% ── SECTION 2: THE CASE AGAINST SUBSTANCE ONTOLOGY (~2,000 words) ────────────

\section{What Quantum Mechanics Requires: The Case Against Substance Ontology}
\label{sec:substance}

\subsection{Substance Ontology and Its Requirements}

Western metaphysics from Aristotle~\citeyearpar{Aristotle_Metaphysics} through
Newton assumed that reality consists of persistent things --- substances ---
that have intrinsic properties and undergo change.
The substance is primary; its relations and processes are secondary.
Things exist first; then they relate.

For substance ontology to be compatible with physics, substances must have
\emph{intrinsic properties}: properties they possess independently of their
relations to other things.
A billiard ball has mass, charge, and position.
These are properties of the ball, not of its relations.

\subsection{Step 1: Bell's Theorem Eliminates Intrinsic Properties}

\citet{Bell1964} proved that no theory assigning definite values to
observables prior to measurement --- no matter how complex or hidden ---
can reproduce the statistical predictions of quantum mechanics without
violating locality.
Experimental tests~\citep{Aspect1982,Zeilinger2023Nobel} have confirmed
Bell inequality violations with loophole-free precision.
Quantum particles do not have determinate spin, position, or momentum
before measurement --- not because we fail to measure them, but because
those properties do not exist prior to the relational event of measurement.

The force of Bell's result deserves emphasis.
It is not a claim about our ignorance.
It is not a claim that there are hidden properties we happen not to know.
Bell's theorem, in conjunction with the experimental record, establishes
that there are no such properties to be ignorant of.
The \citet{Einstein1935EPR} program --- the hypothesis that quantum mechanics
is incomplete and that definite values exist prior to measurement, awaiting
only a more complete theory to reveal them --- is formally closed.
Quantum particles do not have intrinsic properties in the classical sense.
What they have is something different: a relational potential, which is
actualized only through interaction.

\subsection{Step 2: Formal Incompatibility with Substance Ontology}

The consequence for substance ontology follows directly.
A substance, in the classical philosophical sense, is an entity that
possesses intrinsic properties --- properties it holds independently of
its relations to other things.
Bell's theorem has proven that quantum particles do not possess such
properties.
A quantum particle, considered in isolation from relational context, has
no definite position, spin, or momentum.
It is, in Aristotle's framework, not a substance at all.

This is not a matter of quantum mechanics being strange or counterintuitive.
It is a formal incompatibility.
Substance ontology, at its core, requires intrinsic properties.
Quantum mechanics, as confirmed by sixty years of experimental tests of
Bell's inequalities, prohibits them at the most fundamental level of
physical description we possess.
The conclusion is inescapable: substance ontology is not an alternative
interpretation of quantum reality.
It is an interpretation that quantum reality has formally refuted.

The point is worth stating with precision.
The argument is not that substance ontology is philosophically uninteresting
or that process ontology is aesthetically preferable.
The argument is logical: \emph{substance without intrinsic properties
is not substance in any recognizable sense}.
It is an empty placeholder --- a grammatical subject without predicates.
If the only entities at the bottom of physical reality are ones that possess
no intrinsic properties, then the metaphysical category of substance,
understood classically, has no instantiation in the physical world.

\subsection{Step 3: The Wave Function as Description of Potential Events}

If quantum particles are not substances, what are they?
The formalism itself suggests an answer.

The wave function $\psi$ is not a description of what a particle \emph{is}.
It is a description of what relational events are \emph{possible} for it ---
a complete specification of the probabilities of various outcomes, conditional
on various measurement interactions.
The wave function does not say: the particle is here, spinning this way.
It says: if this system interacts with that system in that manner, the
following outcomes are possible with the following probabilities.

This is not a deficiency of the formalism.
It is a feature.
The wave function is precisely calibrated to describe relational potential:
the set of possible events that can actualize when a quantum system enters
a measurement relationship.
Heisenberg's uncertainty principle~\citeyearpar{Heisenberg1927} is not,
on this reading, a limitation on measurement precision.
It is a statement about the ontological structure of quantum systems:
the properties that can be jointly definite are constrained, because
definiteness is not intrinsic but relational, and different measurement
relationships actualize different aspects of the system's potential.

In the vocabulary of process philosophy, the wave function describes an
actual occasion in its \emph{potential phase} --- a concrete quantum system
in the process of becoming, prior to the relational event that will actualize
a definite outcome.
This is not an analogy.
It is a precise ontological mapping.
The wave function is the formal description of what process philosophy would
predict a pre-actualization quantum state to look like: a distribution of
relational possibilities, not a catalogue of intrinsic properties.

\subsection{Step 4: Relational Quantum Mechanics}

\citet{Rovelli1996} makes the process structure of quantum mechanics explicit:
quantum states are always relative to an observer or system --- never absolute.
The measurement event is the moment when superposition resolves into a
definite outcome through relational interaction.
In process ontological language: an actual occasion, a concrete event of
relational becoming in which potential is actualized through interaction.

There is no view from nowhere in Rovelli's relational quantum mechanics.
A particle does not have spin-up; it has spin-up \emph{relative to this
detector, in this measurement event}.
This is not subjectivism --- the outcome is real, and it is determinate.
But it is irreducibly relational: the definiteness of the outcome exists
only within the context of a specific relational event.
Two observers who have not compared notes may assign different states to
the same system --- not because one of them is wrong, but because quantum
states are facts about relations, and different observers stand in
different relational contexts~\citep{Rovelli2021}.

This is precisely the structure that process ontology predicts.
If reality is constituted by relational events rather than intrinsic
properties, then the quantum mechanical description --- in which states
are observer-relative and properties are actualized through interaction ---
is not a puzzle to be explained away.
It is what physics looks like when the underlying metaphysics is correct.

\subsection{Objections and Replies}

The claim that quantum mechanics requires process ontology will be resisted.
Three major objections deserve direct response.

\paragraph{Many-worlds interpretation.}
The Everett many-worlds interpretation~\citep{Penrose1989} avoids the
measurement problem by denying wave function collapse: the wave function
evolves unitarily, and what we experience as a measurement outcome is
our branch of a universal branching structure.
The objection is that many-worlds preserves a form of substance ontology:
the universal wave function is an absolute, observer-independent object.

The reply has two parts.
First, many-worlds resolves the measurement problem by proliferating
ontology on an extraordinary scale: every quantum interaction ramifies
into a new branch of the universal wave function, yielding an uncountably
infinite set of simultaneously real worlds.
This does not so much solve the problem of intrinsic properties as
relocate it: intrinsic properties are now assigned not to particles
but to branches of the multiverse, and the question of what selects
our branch remains unresolved.
Second, many-worlds does not recover what Bell's theorem eliminated.
Even in a many-worlds framework, no single particle in any single
branch has definite intrinsic properties prior to branching.
The relational structure of quantum measurement is preserved in
each branch.
Many-worlds is a response to the \emph{measurement problem}, not
a restoration of \emph{intrinsic properties}.
Bell's argument stands.

\paragraph{QBism.}
Quantum Bayesianism (QBism)~\citep{Stapp2007} treats quantum states as
agents' beliefs about future experiences, not descriptions of objective
reality.
The objection is that if quantum states are just beliefs, the ontological
question does not arise: there is no fact of the matter about what the
particle really is.

The reply is that QBism sidesteps rather than answers the ontological question.
Even if we grant that quantum states are agents' credences, the question
remains: what is the structure of the world such that these credences are
the appropriate ones?
What makes a Bayesian update rational in a quantum context?
QBism can answer these questions only by importing a background ontology ---
and that ontology, if it is to be consistent with quantum mechanics,
will be relational rather than substantival.
QBism defers the ontological question without eliminating it.

\paragraph{Epistemic interpretations.}
A family of epistemic interpretations holds that quantum states describe
not the state of a system but our knowledge of it.
The wave function is a probability distribution over some underlying
classical reality that we are ignorant of.

This is the position that Bell's theorem directly refutes.
The \citet{Einstein1935EPR} argument assumed that there is an underlying
classical reality, with definite properties, that quantum mechanics
incompletely describes.
Bell's theorem proved that no such reality, if it is local, can reproduce
the statistical predictions of quantum mechanics.
Epistemic interpretations are not a live option in the presence of
experimental Bell inequality violations.
The correlations that quantum mechanics predicts --- and that experiment
has confirmed --- cannot be explained by ignorance of pre-existing classical
properties.
They require that the properties themselves come into being through
relational interaction.

The conclusion of all four steps stands: quantum mechanics formally requires
process ontology.
The question is no longer whether to prefer process ontology as an
interpretation.
The question is what process ontology actually says --- and whether any
tradition has previously articulated its structure.
